%%% BEGIN fig:summary-chapter-dag
% Generated by test_ff/scripts/chapter_dag_fig.py from the \uses lists of FINAL_OUTLINE_v2 (edges, badges, counts and
% the users of the tool appendices are data; positions are fixed in the script).  Re-run the script after the verify patch.
\begin{figure}[htbp]
\centering
\begin{tikzpicture}[fd,
  ff ch/.style={draw=black!55,fill=gray!6,line width=0.5pt,rounded corners=1.5pt,minimum width=2.2cm,minimum height=0.42cm,inner sep=0pt},
  ff crit/.style={draw=red!65!black,line width=1.0pt,fill=red!4},
  ff chlab/.style={font=\scriptsize,anchor=west,inner sep=0pt},
  ff badge/.style={circle,inner sep=0pt,minimum size=0.235cm,font=\fontsize{5}{5}\selectfont\bfseries\sffamily,text=white},
  ff badge M/.style={ff badge,fill=OliveGreen!80!black},
  ff badge N/.style={ff badge,fill=RoyalBlue!85!black},
  ff badge O/.style={ff badge,fill=BurntOrange!90!black},
  ff dep/.style={draw=black!45,line width=0.55pt,-{Stealth[length=3.2pt,width=2.8pt]}},
  ff path/.style={draw=red!65!black,line width=0.9pt,-{Stealth[length=4pt,width=3.6pt]}},
  ff pbox/.style={draw=red!65!black,fill=red!4,line width=0.7pt,rounded corners=1.5pt,font=\scriptsize,inner sep=2pt,minimum height=0.42cm},
  ff txt/.style={font=\scriptsize,align=left,anchor=north west,inner sep=0pt,text width=5.0cm},
  ff head/.style={font=\footnotesize\bfseries,anchor=north west,inner sep=0pt}]
  % ---- generated by test_ff/scripts/chapter_dag_fig.py from the \uses data of FINAL_OUTLINE_v2 ----
  \node[ff ch] (Ch01) at (4.20,-0.000) {};
  \node[ff chlab] at ($(Ch01.west)+(0.07,0)$) {\textbf{1}\,引言};
  \node[ff ch] (Ch02) at (6.60,-0.000) {};
  \node[ff chlab] at ($(Ch02.west)+(0.07,0)$) {\textbf{2}\,QCD 基础};
  \node[ff ch] (Ch03) at (5.40,-0.660) {};
  \node[ff chlab] at ($(Ch03.west)+(0.07,0)$) {\textbf{3}\,工具};
  \node[ff ch] (Ch04) at (5.40,-1.320) {};
  \node[ff chlab] at ($(Ch04.west)+(0.07,0)$) {\textbf{4}\,红外结构};
  \node[ff ch] (Ch05) at (5.40,-1.980) {};
  \node[ff chlab] at ($(Ch05.west)+(0.07,0)$) {\textbf{5}\,Sudakov};
  \node[ff ch] (Ch06) at (5.40,-2.640) {};
  \node[ff chlab] at ($(Ch06.west)+(0.07,0)$) {\textbf{6}\,$e^+e^-$};
  \node[ff ch] (Ch07) at (5.40,-3.300) {};
  \node[ff chlab] at ($(Ch07.west)+(0.07,0)$) {\textbf{7}\,DIS/OPE};
  \node[ff badge M] at ($(Ch07.east)+(-0.15,0)$) {M};
  \node[ff ch] (Ch08) at (5.40,-3.960) {};
  \node[ff chlab] at ($(Ch08.west)+(0.07,0)$) {\textbf{8}\,PDF/FF};
  \node[ff badge N] at ($(Ch08.east)+(-0.15,0)$) {N};
  \node[ff ch] (Ch09) at (5.40,-4.620) {};
  \node[ff chlab] at ($(Ch09.west)+(0.07,0)$) {\textbf{9}\,一般方法};
  \node[ff badge O] at ($(Ch09.east)+(-0.15,0)$) {O};
  \node[ff badge N] at ($(Ch09.east)+(-0.40,0)$) {N};
  \node[ff ch] (Ch10) at (5.40,-5.280) {};
  \node[ff chlab] at ($(Ch10.west)+(0.07,0)$) {\textbf{10}\,DIS 全阶};
  \node[ff badge N] at ($(Ch10.east)+(-0.15,0)$) {N};
  \node[ff ch] (Ch11) at (5.40,-5.940) {};
  \node[ff chlab] at ($(Ch11.west)+(0.07,0)$) {\textbf{11}\,共线过程};
  \node[ff badge O] at ($(Ch11.east)+(-0.15,0)$) {O};
  \node[ff badge N] at ($(Ch11.east)+(-0.40,0)$) {N};
  \node[ff ch,ff crit] (Ch12) at (5.40,-6.600) {};
  \node[ff chlab] at ($(Ch12.west)+(0.07,0)$) {\textbf{12}\,DY 过程};
  \node[ff badge O] at ($(Ch12.east)+(-0.15,0)$) {O};
  \node[ff badge N] at ($(Ch12.east)+(-0.40,0)$) {N};
  \node[ff ch] (Ch13) at (4.20,-7.260) {};
  \node[ff chlab] at ($(Ch13.west)+(0.07,0)$) {\textbf{13}\,阈值};
  \node[ff ch,ff crit] (Ch14) at (6.70,-7.260) {};
  \node[ff chlab] at ($(Ch14.west)+(0.07,0)$) {\textbf{14}\,TMD};
  \node[ff badge O] at ($(Ch14.east)+(-0.15,0)$) {O};
  \node[ff badge N] at ($(Ch14.east)+(-0.40,0)$) {N};
  \node[ff ch] (Ch15) at (6.80,-7.920) {};
  \node[ff chlab] at ($(Ch15.west)+(0.07,0)$) {\textbf{15}\,SCET};
  \node[ff ch] (Ch16) at (8.85,-8.580) {};
  \node[ff chlab] at ($(Ch16.west)+(0.07,0)$) {\textbf{16}\,喷注};
  \node[ff ch] (Ch17) at (7.90,-4.620) {};
  \node[ff chlab] at ($(Ch17.west)+(0.07,0)$) {\textbf{17}\,小 $x$};
  \node[ff ch,ff crit] (Ch18) at (8.85,-9.240) {};
  \node[ff chlab] at ($(Ch18.west)+(0.07,0)$) {\textbf{18}\,破坏};
  \node[ff badge O] at ($(Ch18.east)+(-0.15,0)$) {O};
  \node[ff ch] (Ch19) at (1.15,-7.920) {};
  \node[ff chlab] at ($(Ch19.west)+(0.07,0)$) {\textbf{19}\,遍举};
  \node[ff badge N] at ($(Ch19.east)+(-0.15,0)$) {N};
  \node[ff ch] (Ch20) at (6.45,-8.580) {};
  \node[ff chlab] at ($(Ch20.west)+(0.07,0)$) {\textbf{20}\,重夸克};
  \node[ff ch] (Ch21) at (4.40,-7.920) {};
  \node[ff chlab] at ($(Ch21.west)+(0.07,0)$) {\textbf{21}\,幂次修正};
  \node[ff ch,ff crit] (Ch22) at (2.65,-8.580) {};
  \node[ff chlab] at ($(Ch22.west)+(0.07,0)$) {\textbf{22}\,LaMET};
  \node[ff ch,ff crit] (Ch23) at (2.65,-9.240) {};
  \node[ff chlab] at ($(Ch23.west)+(0.07,0)$) {\textbf{23}\,Y.~Liu};
  \node[ff ch,ff crit] (Ch24) at (5.80,-9.900) {};
  \node[ff chlab] at ($(Ch24.west)+(0.07,0)$) {\textbf{24}\,其他前沿};
  \begin{scope}[on background layer]
  \draw[ff dep] (Ch01) -- (Ch03);
  \draw[ff dep] (Ch02) -- (Ch03);
  \draw[ff dep] (Ch03) -- (Ch04);
  \draw[ff dep] (Ch04) -- (Ch05);
  \draw[ff dep] (Ch05) -- (Ch06);
  \draw[ff dep] (Ch06) -- (Ch07);
  \draw[ff dep] (Ch07) -- (Ch08);
  \draw[ff dep] (Ch08) -- (Ch09);
  \draw[ff dep] (Ch08) -- (Ch17);
  \draw[ff dep] (Ch09) -- (Ch10);
  \draw[ff dep,rounded corners=3pt] (Ch09.west) -- (1.15,-4.620) -- (Ch19.north);
  \draw[ff dep] (Ch10) -- (Ch11);
  \draw[ff dep] (Ch11) -- (Ch12);
  \draw[ff dep] (Ch12) -- (Ch13);
  \draw[ff dep] (Ch12) -- (Ch14);
  \draw[ff dep] (Ch13) -- (Ch21);
  \draw[ff dep] (Ch14) -- (Ch15);
  \draw[ff dep] (Ch14) -- (Ch21);
  \draw[ff dep] (Ch15) -- (Ch16);
  \draw[ff dep] (Ch15) -- (Ch20);
  \draw[ff dep] (Ch16) -- (Ch18);
  \draw[ff dep,rounded corners=3pt] (Ch17.east) -- (10.15,-4.620) -- (10.15,-9.900) -- (Ch24.east);
  \draw[ff dep] (Ch18) -- (Ch24);
  \draw[ff dep] (Ch19) -- (Ch20);
  \draw[ff dep] (Ch19) -- (Ch22);
  \draw[ff dep] (Ch20) -- (Ch24);
  \draw[ff dep] (Ch21) -- (Ch22);
  \draw[ff dep] (Ch21) -- (Ch24);
  \draw[ff dep] (Ch22) -- (Ch23);
  \end{scope}
  % ----------------------------------------------------------------- right column: legend, appendices, critical path
  \begin{scope}[xshift=10.42cm]
    \node[ff head] (h1) at (0,0.3) {图例};
    \node[ff ch,ff crit,minimum width=0.62cm,anchor=west] at (0,-0.24) {};
    \node[font=\scriptsize,anchor=west,inner sep=0pt] at (0.75,-0.24) {关键路径上的章（见右下）};
    \node[ff badge M] at (0.12,-0.66) {M}; \node[ff badge N] at (0.4,-0.66) {N}; \node[ff badge O] at (0.68,-0.66) {O};
    \node[font=\scriptsize,anchor=west,inner sep=0pt] at (0.9,-0.66) {该章直接使用附录 M/N/O};
    \draw[ff dep] (0,-1.06) -- (0.62,-1.06);
    \node[font=\scriptsize,anchor=west,inner sep=0pt] at (0.75,-1.06) {依赖（29 条，传递约化）};
    \node[ff head] (h2) at (0,-1.55) {证明附录的使用者};
    \node[ff txt] (t2) at ($(h2.south west)+(0,-0.14)$) {%
      \textbf{M}\,OPE 证明：7（14）\\
      \textbf{N}\,Collins 减除：8（1）、9（23）、10（23）、11（1）、12（3）、14（6）、19（7）\\
      \textbf{O}\,Glauber 相消：9（9）、11（1）、12（27）、14（6）、18（2）\\
      附录之间：M$\to$N（1）、M$\to$O（1）、N$\to$E（1）、N$\to$O（14）、O$\to$N（6）\\
      （括号内为 \texttt{\string\uses} 条数）};
    \node[ff head] (h3) at ($(t2.south west)+(0,-0.3)$) {工具附录的直接使用者};
    \node[ff txt] (t3) at ($(h3.south west)+(0,-0.14)$) {%
      A：3、7、13、19；B：3、5、23、F、L；C：3、6、10、12；D：4、24；E：4；\\
      F：9；I：7、11；J：3、24、N、O；K：3、7、8、10、24、L；L：8、24；G、H、Q 无依赖关系};
    \node[ff head] (h4) at ($(t3.south west)+(0,-0.32)$) {关键路径};
    \node[ff pbox,anchor=north west] (c1) at ($(h4.south west)+(0,-0.16)$) {附录~N};
    \node[ff pbox,anchor=west] (c2) at ($(c1.east)+(0.3,0)$) {附录~O};
    \node[ff pbox,anchor=west] (c3) at ($(c2.east)+(0.3,0)$) {第 12、14、18 章};
    \node[ff pbox,anchor=north east] (c4) at ($(c3.south east)+(0,-0.42)$) {第 22--24 章};
    \node[ff pbox,anchor=east] (c5) at ($(c4.west)+(-0.3,0)$) {第 25 章};
    \node[ff pbox,anchor=east] (c6) at ($(c5.west)+(-0.3,0)$) {附录~P};
    \draw[ff path] (c1) -- (c2); \draw[ff path] (c2) -- (c3);
    \draw[ff path] (c3.south) -- (c3.south |- c4.north);
    \draw[ff path] (c4) -- (c5); \draw[ff path] (c5) -- (c6);
    \node[ff txt] at ($(c6.south -| h4.west)+(0,-0.22)$) {第 25 章与附录~P 汇总全书，自身不含 \texttt{\string\uses}；陈述层次的依赖图见\cref{fig:appP-dependency}。};
  \end{scope}
\end{tikzpicture}
\caption{各章与附录之间的依赖（由 \texttt{\string\uses} 数据汇总）。
左：第~1--24 章，自上而下按最长依赖链分层；箭头 $U\to V$ 表示第 $V$ 章的某个陈述直接或只经由附录陈述使用第 $U$ 章的陈述，图中只画这一关系的传递约化（可由其他箭头推出的箭头略去；第 9 章到第 19 章、第 17 章到第 24 章的两条长箭头绕行）。
章框内的圆标表示该章直接使用证明附录 M、N、O 的陈述（右栏为条数）；附录 M、N、O 本身又使用其前各章的定义（未画出）。
红框与右下角的链为写作与核查的关键路径：附录~N $\to$ 附录~O $\to$ 第 12、14、18 章 $\to$ 第 22--24 章 $\to$ 第 25 章 $\to$ 附录~P。}
\label{fig:summary-chapter-dag}
\end{figure}
%%% END fig:summary-chapter-dag
